\chapter{Tareas Cognitivas y Metricas Conductuales}
\label{cap:cognitive_tasks}

\section{Introduccion}

Ademas de las mediciones fisiologicas continuas de los 4 wearables, el protocolo
FatigueSet incluye \textbf{tareas cognitivas estandarizadas} que miden desempenio
cognitivo y fatiga mental subjetiva. Estas tareas generan 5 archivos CSV por sesion.

\section{Choice Reaction Time (CRT) - Tarea de Tiempo de Reaccion}

\subsection{Objetivo y Protocolo}

Mide velocidad y precision de respuesta motora ante estimulos visuales. Los participantes
ven un circulo en la pantalla y deben presionar un boton lo mas rapido posible.

\begin{itemize}
    \item \textbf{Duracion}: ~2 minutos por bloque
    \item \textbf{Estimulos}: Circulos aleatorios en pantalla
    \item \textbf{Repuesta esperada}: Presion de tecla especifica
    \item \textbf{Metricas}: Tiempo de reaccion, errores
\end{itemize}

\subsection{Archivo: exp\_crt.csv}

Contiene: timestamp de estimulo, tiempo de reaccion (ms), respuesta correcta (0/1)

Filas: ~60-100 registros por bloque

\section{N-Back Task - Prueba de Memoria de Trabajo}

\subsection{Objetivo y Protocolo}

Tarea cognitiva exigente que requiere comparar el estimulo actual con el ''n'' posiciones
anteriores. "1-Back": comparar con el anterior. "2-Back": comparar con hace 2.

\begin{itemize}
    \item \textbf{Duracion}: ~3-5 minutos
    \item \textbf{Dificultad}: 1-Back, 2-Back, 3-Back (niveles crecientes)
    \item \textbf{Estimulos}: Letras (A-Z) mostradas en cuadricula
    \item \textbf{Metricas}: Precision (%), errores, tiempo de reaccion
\end{itemize}

\subsection{Archivo: exp\_nback.csv}

Columnas: stimulus, correctResponse, participantResponse, responseTime (ms), isCorrectResponse

\section{Task Switching}

\subsection{Objetivo y Protocolo}

Mide flexibilidad cognitiva. Participantes ven estimulos (ej: "G6" o "U6") y deben
aplicar una regla (ej: clasificar por letra o numero segun senial).

Cambios de regla generan "switch cost" - slower response times - indicador de
carga cognitiva.

\subsection{Archivo: exp\_task\_switch.csv}

Columnas: stimulus, expectedRule, participantResponse, responseTime, isCorrect

\section{Escala de Fatiga Subjetiva (exp\_fatigue.csv)}

\subsection{Protocolo}

Dentro de cada fase (M1, M2, M3), los participantes completaban dos escalas visuales
analogas (0-100):

\begin{enumerate}
    \item ``Me siento fisicamente cansado'' - Fatiga Fisica
    \item ''Me siento mentalmente cansado'' - Fatiga Mental
\end{enumerate}

\subsection{Archivo: exp\_fatigue.csv}

Columnas: measurementNumber, physicalFatigueScore, mentalFatigueScore (0-100)

Filas: 3 mediciones por sesion (una en cada fase)

\subsection{Interpretacion}

Los autoreportes son subjetivos pero importantes para validar cambios en variables
objetivas. Se espera:

\begin{itemize}
    \item M1 (Baseline): Puntuaciones bajas (10-30)
    \item M2 (Post-ejercicio): Fatiga fisica alta (30-80), mental baja
    \item M3 (Post-mental): Ambas pueden estar elevadas (20-70)
\end{itemize}

\section{Marcadores Temporales (exp\_markers.csv)}

\subsection{Proposito}

Archivo de sincronizacion temporal que marca eventos clave (inicio/fin de tarea,
inicio de fase) para alinear datos de diferentes sensores que pueden tener
desfases temporales.

\subsection{Columnas}

timestamp, event\_name (ej: "START\_EXERCISE", "END\_N\_BACK")

\section{Resumen de Archivos Conductuales}

\begin{table}[H]
\centering
\caption{Resumen de los 5 tipos de archivos conductuales por sesion}
\begin{tabular}{|l|l|l|l|}
\hline
\textbf{Archivo} & \textbf{Registros Aprox.} & \textbf{Frecuencia} & \textbf{Proposito} \\
\hline
exp\_crt.csv & 60-100 & ~1 Hz (eventos) & Velocidad reaccion \\
\hline
exp\_nback.csv & 40-80 & ~1 Hz & Memoria trabajo, precision \\
\hline
exp\_task\_switch.csv & 50-100 & ~1 Hz & Flexibilidad cognitiva \\
\hline
exp\_fatiga.csv & 3 & Discreta (3 fases) & Autoreporte subjetivo \\
\hline
exp\_markers.csv & ~20 & Discreta (eventos) & Sincronizacion temporal \\
\hline
\end{tabular}
\end{table}

Estos archivos conductuales complementan las senales fisiologicas continuas,
permitiendo correlacionar desempenio cognitivo con cambios en HR, EEG, EDA, etc.
